\section{Methods}

\subsection{Implementation, reproducibility, and statistical treatment}

Simulations were implemented in Python using NumPy and PyTorch and run on CPU.
All reported experiments used the same 20 root seeds (51001--51020).
Conditions and plasticity rules within an experiment were paired by root seed.
The experiment scripts write a resolved JSON configuration, compressed NumPy
arrays, and a tidy CSV source-data table; plotting functions generate the
individual panels from those saved arrays. The composed vector figures do not
alter the underlying values.

During preliminary development, bounded hyperparameter configurations were
explored using the covariance matrix adaptation evolution strategy (CMA-ES).
This outer-loop engineering search used reconstruction- or recall-based
objectives to locate workable regimes for population sparsity, activation gain,
EC--CA3 convergence, and plasticity rate. The reported configurations were then
fixed before the 20-seed comparisons. CMA-ES was not applied separately to
individual seeds, experimental conditions, or plasticity rules; it is distinct
from both Adam-based autoencoder pretraining and the local CA3--CA1 updates
studied here.
In generic form, the search proposed bounded parameter vectors $\theta$ and
maximized a development-set score
\begin{equation}
J(\theta)=\frac{1}{N_{\mathrm{dev}}}
 \sum_{n=1}^{N_{\mathrm{dev}}}
 \operatorname{cos}\!\left(
 \widehat x_{\mathrm{EC}}^{(n)}(\theta),x_{\mathrm{EC}}^{(n)}\right),
\end{equation}
or its reconstruction-loss analogue. This procedure selected a computational
operating point; it was not included as evidence for learning, remapping, or
robustness, and optimization generations were not counted as experimental
replicates.

The unit of computational replication was the root seed, not an individual
memory, CA1 unit, track position, or corruption mask. Values were first
averaged over observations nested within a seed and then summarized across
seeds. For a seed-level statistic $y_s$, $s=1,\ldots,S$, curves and error bars
show $\bar y\pm\operatorname{SEM}(y)$, where
\begin{equation}
\bar y=\frac{1}{S}\sum_{s=1}^{S}y_s,\qquad
\operatorname{SEM}(y)=
\sqrt{\frac{1}{S(S-1)}\sum_{s=1}^{S}(y_s-\bar y)^2}.
\end{equation}
Confidence intervals (CIs)
reported in the text are two-sided 95\% Student-$t$ intervals over the 20
seed-level values,
$\bar y\pm t_{0.975,S-1}\operatorname{SEM}(y)$. Contrasts use paired
seed-level differences $d_s=y_s^{(A)}-y_s^{(B)}$ rather than differences
between independently aggregated conditions. These intervals
quantify sensitivity to sampled initializations and inputs, not uncertainty
over a biological population. No null-hypothesis significance tests or
multiple-comparison corrections were applied; all analyses are descriptive
simulation results.

\subsection{EC patterns and the pretrained CA1 basis}

EC input, CA1, and CA3-like populations each contained 50 units. A one-hidden-
layer autoencoder mapped EC input to a 50-unit CA1 representation and back to a
50-unit EC output. CA1 activity used a differentiable top-$K$ sparse activation
with $K_{\mathrm{CA1}}=5$ and inverse temperature $\beta=25$; output activity
used a sigmoid decoder with gain 25. Linear encoder and decoder biases were
disabled. Specifically, if $z_{(1)}\geq\cdots\geq z_{(N)}$ are the sorted
preactivations, the activation threshold and sparse nonlinearity were
\begin{align}
\vartheta_K(z) &= \frac{z_{(K)}+z_{(K+1)}}{2},\\
[\mathcal{S}_{K,\beta}(z)]_i
 &= \sigma\!\left(\beta[z_i-\vartheta_K(z)]\right),
\qquad \sigma(u)=\frac{1}{1+e^{-u}}.
\label{eq:sparse-activation}
\end{align}
Thus approximately $K$ units lay above the adaptive half-activation threshold,
while gradients remained defined near that threshold. For EC pattern
$x_{\mathrm{EC}}\in\{0,1\}^{50}$, encoder $E$, and decoder $D$, pretraining
computed
\begin{align}
x_{\mathrm{IS}} &= \mathcal{S}_{5,25}(E x_{\mathrm{EC}}),\\
\widetilde{x}_{\mathrm{EC}} &= \sigma(25D x_{\mathrm{IS}}),\\
\mathcal{L}_{\mathrm{AE}} &=
 \frac{1}{N}\sum_{n=1}^{N}
 \left\|\widetilde{x}_{\mathrm{EC}}^{(n)}
             -x_{\mathrm{EC}}^{(n)}\right\|_2^2.
\label{eq:autoencoder}
\end{align}
The weights were trained with Adam (learning rate 0.001, batch size 64) to
minimize Eq.~\ref{eq:autoencoder}. Pretraining occurred before
CA3--CA1 storage, after which the encoder and decoder were fixed. The encoder
therefore supplied the instructed CA1 state (IS), and the decoder defined the
output coordinates against which compatibility was tested.

For the readout-compatibility experiment, the autoencoder was trained for 1024
epochs on 2048 unique binary EC patterns with five active units and evaluated
on 256 held-out patterns. Twenty-eight additional unique patterns were reserved
for one-shot associative storage. For the cue-reconfiguration experiment, the
autoencoder was trained for 256 epochs on the 40 training laps and evaluated on
10 separately generated laps. For the degradation experiment, it was trained
for 256 epochs on 20 intact laps and evaluated on five separately generated
laps. Each lap contained 50 positions.

Track inputs contained 25 MEC-like and 25 LEC-like units. For position
$p\in\{0,\ldots,L-1\}$ on a circular track of length $L=50$, circular distance
was $d_L(p,\mu)=\min(|p-\mu|,L-|p-\mu|)$. MEC unit $i$ had a preferred position
$\mu_i$ from an even tiling of the track and Gaussian activation probability
\begin{equation}
g_i(p)=\exp\!\left[-\frac{1}{2}
 \left(\frac{d_L(p,\mu_i)}{\sigma_{\mathrm{MEC}}}\right)^2\right],
\qquad \sigma_{\mathrm{MEC}}=5.
\end{equation}
Binary MEC activity was sampled independently as
$x_{\mathrm{MEC},i}(p)\sim\operatorname{Bernoulli}(g_i(p))$. Two fixed,
non-overlapping 25-dimensional LEC cue templates $q_1$ and $q_2$ were centered
at $c_1=10$ and $c_2=30$. Cue $j$ was present with probability
\begin{align}
g_j^{\mathrm{cue}}(p)
 &=\exp\!\left[-\frac{1}{2}
 \left(\frac{d_L(p,c_j)}{4}\right)^2\right],\\
P_j(p)&=\operatorname{clip}_{[0,1]}
 \left\{2\sigma\!\left(40[g_j^{\mathrm{cue}}(p)-0.1]\right)-1\right\}.
\end{align}
If the Bernoulli cue-presence sample was positive, the LEC block was set to
$q_j$; otherwise it was zero. The complete entorhinal target was the
concatenation $x_{\mathrm{EC}}(p)=[x_{\mathrm{MEC}}(p);
x_{\mathrm{LEC}}(p)]$. This stochastic construction generated different laps
from the same underlying spatial and cue statistics.

\subsection{CA3-like retrieval keys and CA3--CA1 plasticity}

EC input was projected to 50 CA3-like units through a fixed sparse matrix
followed by a differentiable top-$K$ activation. The compatibility experiment
used two EC inputs per CA3-like unit, $K_{\mathrm{CA3}}=5$, inverse-temperature
parameter $\beta_{\mathrm{CA3}}=200$, and learning rate $\alpha=0.08$. Track
experiments used 29 EC inputs per CA3-like unit, $K_{\mathrm{CA3}}=3$,
$\beta_{\mathrm{CA3}}=170.6$, and $\alpha=0.092$. The normal track map was
constructed to balance EC participation across CA3-like units.
CA3--CA1 weights were initialized to zero. The CA1 recall activation used
$K_{\mathrm{CA1}}=5$, with inverse temperature 25 in the compatibility
experiment and 41.8 in the track experiments.

Writing and recall used the same fixed EC--CA3 projection $A$ and frozen
decoder $D$. At weight state $W_t$, the complete forward computation was
\begin{align}
x_{\mathrm{CA3}} &=
 \mathcal{S}_{K_{\mathrm{CA3}},\beta_{\mathrm{CA3}}}
 (A x_{\mathrm{EC}}),\\
x_{\mathrm{CA1}} &=
 \mathcal{S}_{5,\beta_{\mathrm{CA1}}}(W_t x_{\mathrm{CA3}}),\\
\widehat{x}_{\mathrm{EC}} &= \sigma(25D x_{\mathrm{CA1}}).
\label{eq:recall-forward}
\end{align}
Each selected input to a CA3-like unit had projection weight $1/50$; all other
entries of $A$ were zero. Consequently, the convergence parameter controlled
which EC components contributed to a key, whereas the top-$K$ nonlinearity
controlled key population sparsity. Equation~\ref{eq:recall-forward} was also
used to compute the pre-update recall state required by the error-driven rule.

Let $x_{\mathrm{EC}}$ denote the EC input/target,
$x_{\mathrm{CA3}}$ its retrieval key, $x_{\mathrm{IS}}$ the instructed CA1
state produced by the fixed encoder, and $x_{\mathrm{CA1}}$ the state recalled
through the current CA3--CA1 weights $W_t$. The instructive-driven update was
\begin{equation}
W_{t+1} =
  (1-\alpha x_{\mathrm{IS}})\odot W_t
  + \alpha x_{\mathrm{IS}}x_{\mathrm{CA3}}^{\mathsf T},
\label{eq:instructive-driven-rule}
\end{equation}
where the first term is broadcast across presynaptic dimensions. For the
bounded error-driven rule,
\begin{align}
e^+ &= [x_{\mathrm{IS}}-x_{\mathrm{CA1}}]_+, &
e^- &= [x_{\mathrm{CA1}}-x_{\mathrm{IS}}]_+,\\
\delta^+ &= \alpha(e^+x_{\mathrm{CA3}}^{\mathsf T})(1-W_t), &
\delta^- &= \alpha(e^-x_{\mathrm{CA3}}^{\mathsf T})W_t,\\
W_{t+1} &= \operatorname{clip}_{[0,1]}
           (W_t+\delta^+-\delta^-).
\label{eq:error-driven-rule}
\end{align}
Recalled CA1 activity was obtained by applying the same target sparsity to
$W_tx_{\mathrm{CA3}}$ and was passed through the fixed pretrained decoder.
The instructive-driven rule depends on the instructed target and presynaptic key but not
the current recall error. The error-driven rule separates the residual into
positive and negative components, such that under- and overexpressed CA1
components drive bounded potentiation and depression, respectively.
All products between equally shaped arrays were elementwise unless an outer
product or matrix product is shown explicitly. Memories were presented
sequentially, and $W_{t+1}$ was carried to the next item or position without
replay, batching, or gradient descent through the local plasticity dynamics.

\subsection{Paired plasticity-rule design}

Both rules were evaluated in every main experiment. Within a root seed they
shared EC inputs, pretrained autoencoder weights, CA3 wiring seeds, storage
order, and--in the degradation experiment--corruption masks. Each rule had its
own CA3--CA1 weight matrix, which evolved according to the corresponding
update. All memory hyperparameters were shared rather than optimized for each
rule. Supplementary Fig.~\ref{fig:rule-comparison} collects four paired contrasts
derived directly from the main result arrays; it is not a separate simulation.

\subsection{Readout-compatibility experiment}

For each seed, the 28 held-out EC patterns were assigned CA3-like retrieval keys
and pretrained CA1 instructions, then written once in a shared random order.
The aligned condition used the original IS and decoder. The fixed-permutation
condition applied one seed-specific CA1 permutation to every IS while leaving
the decoder unchanged. The matched-decoder condition used the same permuted
instructions and correspondingly permuted decoder columns. The random-content
condition associated each key with a shuffled memory instruction. The
no-plasticity condition retained zero CA3--CA1 weights. Recall was evaluated
using cosine similarity between reconstructed EC output and the corresponding
EC target. Each seed's plotted value is its mean across 28 memories.

More explicitly, let $P$ be the permutation matrix and $h=x_{\mathrm{IS}}$.
The aligned condition stored $h$ and read it with $D$; the fixed-permutation
condition stored $Ph$ but retained $D$; and the matched-decoder condition used
$Ph$ together with $D'=DP^{\mathsf T}$. In the idealized exact-recall case,
\begin{equation}
D'Ph=DP^{\mathsf T}Ph=Dh,
\end{equation}
so matched decoder relabeling restores the original output coordinates without
restoring the original CA1 unit identities. This control isolates coordinate
compatibility from a generic effect of permutation. Similarity between two
nonzero vectors $u$ and $v$ was
\begin{equation}
\operatorname{cos}(u,v)=\frac{u^{\mathsf T}v}
 {\lVert u\rVert_2\lVert v\rVert_2};
\end{equation}
zero-norm cases were assigned similarity zero by the numerical implementation.

\subsection{Cue-reconfiguration experiment}

Models learned 40 complete laps. Cue identities exchanged positions every 10
laps, alternating between assignments $[0,1]$ and $[1,0]$. The matched
no-swap schedule retained $[0,1]$ for all 40 laps. Within a seed, the schedules
shared the pretrained autoencoder, initial CA3--CA1 weights, MEC-like
trajectory, lap count, plasticity parameters, and the exact EC--CA3 projection;
only cue identity after the first 10 laps differed by design. Both plasticity
rules used this same projection. Thus all schedule contrasts in
Fig.~\ref{fig:compatibility}B,E,F are paired at the level of the root seed and
fixed retrieval-key wiring.

After each training lap, plasticity was paused while both fixed cue contexts
were probed; the probe matching the context scheduled for that lap was retained
for the transition analysis. For each CA1 unit, activity was mean-centered
across positions before calculating cosine similarity between consecutive
spatial tuning curves. Similarities were then averaged across units. Cue-exchange
transitions were compared with the same lap boundaries in the paired no-swap
model.

Let $f_{\ell c i}(p)$ be the frozen-probe activity of CA1 unit $i$ at position
$p$, after training lap $\ell$, under cue context $c$, and define the centered
curve $\widetilde f_{\ell c i}(p)=f_{\ell c i}(p)-L^{-1}
\sum_{r=0}^{L-1}f_{\ell c i}(r)$. Consecutive-lap tuning similarity was
\begin{equation}
T_\ell=\frac{1}{N_{\mathrm{CA1}}}\sum_{i=1}^{N_{\mathrm{CA1}}}
 \operatorname{cos}(\widetilde f_{\ell c_\ell i},
                    \widetilde f_{\ell-1,c_{\ell-1},i}).
\label{eq:transition-similarity}
\end{equation}
The transition statistic compared $T_\ell$ at boundaries
$\ell\in\{10,20,30\}$ with the same boundaries in the matched no-swap run.

For the exact-wiring event analysis, each boundary $e\in\{10,20,30\}$ was
treated as relative lap $\tau=0$, the first lap learned under the exchanged
cue assignment. The analysis window was $\tau\in\{-3,\ldots,9\}$. Let
$c_{e+\tau}$ be the context scheduled at that relative lap. The final
pre-exchange probe and the final probe of the ensuing ten-lap block defined
the old and adapted references, respectively:
\begin{align}
R_{\mathrm{old}}(e,\tau)
 &=\frac{1}{N_{\mathrm{CA1}}}\sum_i
 \operatorname{cos}\!\left(
 \widetilde f_{e+\tau,c_{e+\tau},i},
 \widetilde f_{e-1,c_{e-1},i}\right),\\
R_{\mathrm{adapt}}(e,\tau)
 &=\frac{1}{N_{\mathrm{CA1}}}\sum_i
 \operatorname{cos}\!\left(
 \widetilde f_{e+\tau,c_{e+\tau},i},
 \widetilde f_{e+9,c_{e+9},i}\right).
\label{eq:remapping-references}
\end{align}
The adapted reference is retrospective: it describes convergence toward the
state reached at the end of a block and is not a prediction of a future map.
Consecutive-lap similarity used Eq.~\ref{eq:transition-similarity} at each
$e+\tau$. Values were first averaged over the three events within a root seed
and then summarized across seeds. The no-swap control was indexed by the same
nominal boundaries, although its cue assignment did not change.

For the population representational-similarity matrix, the complete
position-by-unit probe map after lap $\ell$ was centered independently for each
unit over position and vectorized as
\begin{equation}
v_\ell=\operatorname{vec}\!\left[
f_\ell-\frac{1}{L}\mathbf{1}\mathbf{1}^{\mathsf T}f_\ell
\right],\qquad
P_{\ell m}=\operatorname{cos}(v_\ell,v_m).
\label{eq:population-map-similarity}
\end{equation}
The displayed matrix is the elementwise mean of $P$ across 20 root seeds.
Unit-level endpoint similarity was
$U_{ei}=\operatorname{cos}(\widetilde f_{e-1,c_{e-1},i},
\widetilde f_{e+9,c_{e+9},i})$. Its empirical cumulative distribution was
computed within each seed after pooling that seed's three exchange events;
the plotted curve and uncertainty were then calculated across seeds, avoiding
treatment of individual CA1 units as independent replicates.

After the final lap, plasticity remained frozen and the two cue contexts were
probed. For each CA1 unit, spatial stability was the cosine similarity between
its two mean-centered context tuning curves. Cue modulation was the mean
absolute activity difference between contexts. These metrics quantify context
sensitivity but do not themselves classify rate or global remapping. Heat maps
show the error-driven seed whose mean stability and modulation were closest to the
across-seed medians; units were ordered by preferred position across the two
contexts.
For final-lap contexts $a$ and $b$, the plotted unit-level quantities were
\begin{align}
S_i &= \operatorname{cos}(\widetilde f_{a i},\widetilde f_{b i}),\\
M_i &= \frac{1}{L}\sum_{p=0}^{L-1}|f_{a i}(p)-f_{b i}(p)|.
\end{align}
Here $S_i$ measures preservation of spatial profile after removing mean firing,
whereas $M_i$ retains absolute context-dependent rate differences; their joint
distribution therefore separates two complementary aspects of mixed tuning.

\subsection{Selective degradation experiments}

Models stored 20 intact laps with cue assignment $[0,1]$. During frozen recall,
a fixed random subset of eligible EC units was zeroed throughout a complete
probe lap. Masks were sampled independently between probes. Dropped fractions
were 0, 0.25, 0.50, 0.75, and 0.90, with 12 masks per fraction and root seed.
LEC-like degradation operated on the final 25 input units; MEC-like degradation
operated on the first 25.

For dropped fraction $q$ and mask replicate $r$, a binary mask
$m^{(q,r)}\in\{0,1\}^{50}$ set
$\operatorname{round}(25q)$ randomly selected coordinates in the targeted
input half to zero and left all other coordinates at one. The same mask was
held fixed for all positions of that probe lap:
\begin{equation}
x_{\mathrm{EC}}^{(q,r)}(p)=m^{(q,r)}\odot x_{\mathrm{EC}}(p).
\label{eq:degradation-mask}
\end{equation}
Plasticity was disabled during these probes, so degradation altered retrieval
evidence but did not rewrite the stored association.

Cue identity was decoded at the two cue positions by cosine comparison between
recalled LEC output and the two fixed cue templates. Spatial accuracy was the
fraction of recalled MEC outputs whose nearest clean MEC target had the correct
track-position index; chance was therefore $1/50$. MEC output cosine was
averaged across all positions. For every curve, the 12 mask-level values were
first averaged within a seed before means and uncertainty were calculated
across seeds.

Writing $\widehat x_{\mathrm{LEC}}(c_j)$ for recalled LEC output at cue center
$c_j$ and $x_{\mathrm{MEC}}^*(p)$ for the clean MEC target at position $p$,
the three reported metrics were
\begin{align}
A_{\mathrm{cue}}
 &=\frac{1}{2}\sum_{j=1}^{2}
 \mathbf{1}\!\left[j=\arg\max_{k\in\{1,2\}}
 \operatorname{cos}(\widehat x_{\mathrm{LEC}}(c_j),q_k)\right],\\
A_{\mathrm{pos}}
 &=\frac{1}{L}\sum_{p=0}^{L-1}
 \mathbf{1}\!\left[p=\arg\max_{u\in\{0,\ldots,L-1\}}
 \operatorname{cos}(\widehat x_{\mathrm{MEC}}(p),x_{\mathrm{MEC}}^*(u))\right],\\
C_{\mathrm{MEC}}
 &=\frac{1}{L}\sum_{p=0}^{L-1}
 \operatorname{cos}(\widehat x_{\mathrm{MEC}}(p),x_{\mathrm{MEC}}^*(p)).
\end{align}
Cue accuracy has two-class chance level $1/2$, whereas exact-position accuracy
has chance level $1/L=1/50$. MEC cosine is graded and has no corresponding
uniform-classification chance line.

The dense common-key control replaced the EC--CA3 map with a matrix whose every
entry was $1/50$. All CA3-like units consequently received the same average EC
input and generated a common, non-item-specific key. Because this control also
changes connection density and sparsity relative to the normal map, it tests
the need for key discriminability rather than the isolated effect of sparse
connectivity.
In matrix notation the control used $A_{ij}=1/50$ for all $i,j$; hence every
CA3 preactivation was identical for a given input. Equation~\ref{eq:sparse-activation}
therefore returned the same $0.5$-valued vector for every input, rather
than an item-specific key.
